\documentclass[12pt]{article}

\usepackage[margin=1in]{geometry}
\usepackage{enumitem}
\usepackage{amsmath}
\usepackage{amssymb}
\usepackage{array}
\usepackage{titlesec}

\title{Lecture 10 Notes:\\ Aristotle's \textit{Topics}, Book III}
\author{}
\date{}

\begin{document}

\maketitle

\section*{Central Theme}

Aristotle's \textit{Topics}, Book III gives commonplace rules for comparing
goods, desirables, ends, means, consequences, and circumstances.

\[
\boxed{
\textit{Topics}, \text{ Book III}
=
\text{commonplace rules for comparing goods and desirables}
}
\]

Book II taught how to test whether a predicate belongs accidentally. Book III
teaches how to compare goods when the question is not simply whether something
is good, but which of two things is better, more desirable, or more worth
choosing.

\section*{1. From Accidental Predication to Comparison}

Book II focused on questions of accidental predication.

\[
\text{Does } P \text{ belong to } S?
\]

Book III shifts to comparative questions.

\[
\text{Is } A \text{ more desirable than } B?
\]

\[
\text{Is } A \text{ better than } B?
\]

\[
\text{Is } A \text{ more objectionable than } B?
\]

\subsection*{Teaching Point}

Book III turns dialectic toward practical comparison and choice.

\[
\boxed{
\text{Dialectic is useful not only for testing claims, but also for weighing alternatives.}
}
\]

\section*{2. The Scope of Comparative Inquiry}

Aristotle says that comparison is most useful when the things compared are
nearly related.

We do not usually need dialectical rules to decide between things that are
wildly different in value.

\subsection*{Example}

No one seriously wonders whether happiness or wealth is more desirable.

\[
\text{happiness} > \text{wealth}
\]

The real work of comparison concerns close cases.

\[
\text{justice vs. courage}
\]

\[
\text{health vs. strength}
\]

\[
\text{prudence vs. courage in different seasons of life}
\]

\subsection*{Teaching Point}

Comparative dialectic is most useful where choice is genuinely contestable.

\[
\boxed{
\text{The closer the alternatives, the more useful the topical rules become.}
}
\]

\section*{3. The Basic Strategy of Book III}

When two goods are closely related and no superiority is immediately obvious,
we look for an advantage on one side.

\[
A \approx B
\]

\[
A + \text{advantage}
\]

\[
\therefore A \text{ is more desirable}
\]

One relevant advantage, or several together, can turn the judgment.

\subsection*{Teaching Point}

Book III gives a list of advantage-tests.

\[
\boxed{
\text{In close comparisons, a single relevant advantage can decide the argument.}
}
\]

\section*{4. More Lasting and More Secure Goods}

The first rule is that what is more lasting or more secure is more desirable
than what is less lasting or less secure.

\[
\text{more lasting} > \text{less lasting}
\]

\[
\text{more secure} > \text{less secure}
\]

\subsection*{Examples}

\[
\text{stable virtue} > \text{temporary pleasure}
\]

\[
\text{secure friendship} > \text{unstable popularity}
\]

\subsection*{Teaching Point}

Durability and security are marks of greater desirability.

\[
\boxed{
\text{A good that lasts or is secure has an advantage over a fragile good.}
}
\]

\section*{5. What the Prudent, Good, Lawful, or Expert Would Choose}

Another rule is to ask what would be chosen by:

\[
\begin{array}{ll}
1. & \text{the prudent person} \\
2. & \text{the good person} \\
3. & \text{the right law} \\
4. & \text{the expert in the relevant field}
\end{array}
\]

\subsection*{Example}

In medicine, the more desirable treatment is the one that most or all doctors
would choose as doctors.

In carpentry, the more desirable tool is the one that most or all carpenters
would choose as carpenters.

\subsection*{Teaching Point}

Choice by the competent is evidence of greater desirability.

\[
\boxed{
\text{The judgment of the competent supplies a standard for comparison.}
}
\]

\section*{6. Absolute and Relative Standards}

Aristotle distinguishes absolute and relative standards of comparison.

\[
\text{absolute standard} = \text{the verdict of the better science}
\]

\[
\text{relative standard} = \text{the science or craft relevant to this person}
\]

\subsection*{Teaching Point}

The standard of comparison depends on whether the question is absolute or
relative.

\[
\boxed{
\text{Better simply and better for this person are not always the same question.}
}
\]

\section*{7. What Belongs to the Genus of Good}

What is itself called a good is more desirable than what does not fall directly
within the genus of good.

\subsection*{Example}

Justice is more desirable than a just man in this respect.

\[
\text{justice} \in \text{good}
\]

A just man is good by possessing justice, but he is not called ``a good'' in the
same direct generic way.

\subsection*{Teaching Point}

What belongs directly to the relevant genus outranks what is only qualified by
it.

\[
\boxed{
\text{Direct membership in the genus of good carries comparative force.}
}
\]

\section*{8. Desired for Itself and Desired for Something Else}

What is desired for itself is more desirable than what is desired for something
else.

\[
\text{for itself} > \text{for another}
\]

\subsection*{Example}

\[
\text{health} > \text{gymnastics}
\]

Health is desired for itself.

Gymnastics is desired for the sake of health.

\subsection*{Teaching Point}

Ends outrank means.

\[
\boxed{
\text{What is chosen for its own sake is more desirable than what is chosen as an instrument.}
}
\]

\section*{9. Desirable in Itself and Desirable Per Accidens}

What is desirable in itself is more desirable than what is desirable merely per
accidens.

\subsection*{Example}

Justice in friends is desirable in itself.

Justice in enemies may be desired only accidentally, because it prevents them
from harming us.

\[
\text{justice in friends} > \text{justice in enemies}
\]

\subsection*{Teaching Point}

A good chosen for its own character outranks a good chosen only because of side
effects.

\[
\boxed{
\text{In itself} > \text{per accidens}
}
\]

\section*{10. Cause of Good in Itself and Cause of Good Per Accidens}

What is in itself the cause of good is more desirable than what is only
accidentally the cause of good.

\subsection*{Example}

\[
\text{virtue} > \text{luck}
\]

Virtue is a cause of good in itself.

Luck is a cause of good only per accidens.

\subsection*{Contrary Case}

Likewise, what is in itself the cause of evil is more objectionable than what is
accidentally the cause of evil.

\[
\text{vice} > \text{chance as a cause of evil}
\]

\subsection*{Teaching Point}

Stable intrinsic causes outrank accidental sources of benefit.

\[
\boxed{
\text{A good cause by nature is better than a lucky cause by accident.}
}
\]

\section*{11. Good Absolutely and Good for a Particular Person}

What is good absolutely is more desirable than what is good only for a
particular person.

\subsection*{Example}

\[
\text{recovery of health} > \text{surgical operation}
\]

Recovery is good absolutely.

Surgery is good only for the person who needs it.

\subsection*{Teaching Point}

Conditional goods are lower than goods desirable without special need.

\[
\boxed{
\text{Good absolutely} > \text{good for this person under this condition}
}
\]

\section*{12. Good by Nature and Acquired Good}

What is good by nature is more desirable than what is not good by nature.

\subsection*{Example}

Justice is better than the just man in this respect.

\[
\text{justice} = \text{good by nature}
\]

\[
\text{just man} = \text{good by acquisition}
\]

\subsection*{Teaching Point}

A natural good has priority over a good that possesses goodness derivatively.

\[
\boxed{
\text{Good by nature} > \text{good by acquisition}
}
\]

\section*{13. Better Subject, Better Attribute}

An attribute that belongs to a better and more honorable subject is more
desirable.

\subsection*{Examples}

\[
\text{what belongs to a god} > \text{what belongs to a man}
\]

\[
\text{what belongs to the soul} > \text{what belongs to the body}
\]

Likewise, the property of a better thing is better than the property of a worse
thing.

\subsection*{Teaching Point}

The rank of a subject can elevate the rank of what belongs to it.

\[
\boxed{
\text{Better subject} \Rightarrow \text{better attribute}
}
\]

\section*{14. What Is Inherent in Better, Prior, or More Honorable Things}

What is inherent in better, prior, or more honorable things is better.

\subsection*{Example}

Health is better than strength or beauty.

\[
\text{health} > \text{strength}
\]

\[
\text{health} > \text{beauty}
\]

Health is rooted in the primary constituents of the animal: the moist, dry, hot,
and cold.

Strength belongs to sinews and bones.

Beauty belongs to symmetry of limbs.

\subsection*{Teaching Point}

A good rooted in more primary parts has greater rank.

\[
\boxed{
\text{The level at which a good inheres affects its rank.}
}
\]

\section*{15. Ends and Means}

The end is generally more desirable than the means.

\[
\text{end} > \text{means}
\]

Of two means, the one nearer the end is more desirable.

\[
\text{nearer means} > \text{more remote means}
\]

\subsection*{Example}

A means to happiness is more desirable than a means to some lesser end.

\[
\text{means to happiness} > \text{means to lesser end}
\]

\subsection*{Teaching Point}

Value increases as we move closer to the final end.

\[
\boxed{
\text{Ends rank means, and final ends rank subordinate ends.}
}
\]

\section*{16. Competence and Incompetence}

Aristotle briefly notes that the competent is more desirable than the
incompetent.

\[
\text{competent} > \text{incompetent}
\]

\subsection*{Teaching Point}

A capacity ordered well toward its function outranks failure of capacity.

\[
\boxed{
\text{Functioning well is better than failing to function.}
}
\]

\section*{17. Productive Agents and Their Ends}

Of two productive agents, the one whose end is better is more desirable.

\[
\text{better end} \Rightarrow \text{better productive agent}
\]

\subsection*{Example}

If one thing produces happiness and another produces health, and happiness
exceeds health more than health exceeds what produces health, then what
produces happiness is more desirable than health.

\subsection*{Teaching Point}

The rank of an end can raise the rank of what produces it.

\[
\boxed{
\text{Compare productive agents by comparing what they produce.}
}
\]

\section*{18. Noble, Precious, and Praiseworthy Things}

What is nobler, more precious, and more praiseworthy in itself is more
desirable.

\subsection*{Examples}

\[
\text{friendship} > \text{wealth}
\]

\[
\text{justice} > \text{strength}
\]

Friendship is prized for itself.

Wealth is prized for something else.

Justice is praiseworthy in itself.

Strength may be prized for what it can accomplish.

\subsection*{Teaching Point}

What is honored for itself outranks what is valued merely as useful.

\[
\boxed{
\text{Intrinsic nobility outranks instrumental usefulness.}
}
\]

\section*{19. Judge by Consequences}

When two things are very much alike and no superiority is visible, examine their
consequences.

The thing followed by the greater good is more desirable.

If the consequences are evil, the thing followed by the lesser evil is more
desirable.

\subsection*{Pattern}

\[
A \rightarrow \text{greater good}
\]

\[
B \rightarrow \text{lesser good}
\]

\[
\therefore A > B
\]

\subsection*{Prior and Later Consequences}

Aristotle distinguishes prior and later consequences.

If a person learns, then:

\[
\text{before: he was ignorant}
\]

\[
\text{after: he knows}
\]

As a rule, the later consequence is better to consider.

\subsection*{Teaching Point}

When the thing itself is hard to compare, compare what follows from it.

\[
\boxed{
\text{Consequences can reveal hidden superiority.}
}
\]

\section*{20. More Goods and Fewer Goods}

A greater number of good things is generally more desirable than a smaller
number.

\[
\text{more goods} > \text{fewer goods}
\]

This is especially clear when the smaller number is included in the greater.

\subsection*{Important Qualification}

If one item is desired only for the sake of another, the two together are not
necessarily better than the one.

\subsection*{Example}

\[
\text{recovery of health} + \text{health}
\]

is not necessarily better than:

\[
\text{health}
\]

because recovery is desired for the sake of health.

\subsection*{Teaching Point}

Counting goods works only when the added item contributes independent value.

\[
\boxed{
\text{Means plus end is not always better than the end alone.}
}
\]

\section*{21. Pleasure and Pain as Modifiers}

The same things are more valuable when accompanied by pleasure.

\[
\text{good} + \text{pleasure} > \text{good without pleasure}
\]

They are also more valuable when free from pain.

\[
\text{good without pain} > \text{good with pain}
\]

\subsection*{Teaching Point}

Pleasure can enhance a good, and pain can reduce its desirability.

\[
\boxed{
\text{Pleasure and pain modify desirability without becoming the whole standard.}
}
\]

\section*{22. Season and Timing}

A good is more desirable at the season when it is of greater consequence.

\subsection*{Examples}

Freedom from pain is more desirable in old age than in youth.

\[
\text{freedom from pain in old age} > \text{freedom from pain in youth}
\]

Prudence is more desirable in old age.

Courage may be more urgently needed in youth.

Temperance is especially needed by the young because they are more troubled by
their passions.

\subsection*{Teaching Point}

The value of a good may rise or fall with the season of life.

\[
\boxed{
\text{Timing can change comparative desirability.}
}
\]

\section*{23. Useful at All or Most Seasons}

What is useful at every season or most seasons is more desirable than what is
useful only sometimes.

\subsection*{Example}

Justice and temperance are more useful than courage in this respect.

\[
\text{justice and temperance} > \text{courage}
\]

Justice and temperance are always useful.

Courage is useful only at certain times.

\subsection*{Teaching Point}

Constant usefulness is a mark of greater desirability.

\[
\boxed{
\text{A good useful across more circumstances has comparative advantage.}
}
\]

\section*{24. If All Possess One Good, Is the Other Still Needed?}

Aristotle gives a powerful comparative rule.

If all possessed one good and no longer needed the other, then the first good is
more desirable.

\subsection*{Example}

If everyone were just, courage would not be needed.

But if everyone were courageous, justice would still be needed.

Therefore:

\[
\text{justice} > \text{courage}
\]

\subsection*{Pattern}

\[
\text{If all have } A, B \text{ becomes unnecessary}
\]

\[
\text{If all have } B, A \text{ remains necessary}
\]

\[
\therefore A > B
\]

\subsection*{Teaching Point}

The good that makes the other unnecessary has priority.

\[
\boxed{
\text{A more fundamental good can render a secondary good unnecessary.}
}
\]

\section*{25. Destructions, Losses, Generations, Acquisitions, and Contraries}

Aristotle says to judge goods by:

\[
\begin{array}{ll}
1. & \text{their destructions} \\
2. & \text{their losses} \\
3. & \text{their generations} \\
4. & \text{their acquisitions} \\
5. & \text{their contraries}
\end{array}
\]

Things whose destruction or loss is more objectionable are themselves more
desirable.

Things whose acquisition or generation is more desirable are themselves more
desirable.

\subsection*{Teaching Point}

The value of a thing appears in the value of gaining it or losing it.

\[
\boxed{
\text{A greater loss indicates a greater good.}
}
\]

\section*{26. Nearness and Likeness to the Good}

What is nearer to the good is more desirable.

\[
\text{nearer to good} > \text{farther from good}
\]

What more resembles a better thing may also be more desirable.

\subsection*{Warning}

Likeness can mislead.

A monkey resembles a human more than a horse does, but that does not make the
monkey more beautiful.

A resemblance may be slight, irrelevant, caricatured, or degrading.

\subsection*{Teaching Point}

Likeness is useful only when it tracks the relevant excellence.

\[
\boxed{
\text{Mere resemblance is not the same as real superiority.}
}
\]

\section*{27. More Conspicuous, More Difficult, More Personal, Freer from Evil}

Aristotle gives several shorter comparative rules.

\[
\text{more conspicuous good} > \text{less conspicuous good}
\]

\[
\text{more difficult good} > \text{easier good}
\]

\[
\text{more personal possession} > \text{more widely shared possession}
\]

\[
\text{freer from evil} > \text{attended by evil}
\]

\subsection*{Teaching Point}

Rarity, difficulty, personal possession, and purity from evil can all increase
desirability.

\[
\boxed{
\text{A good may be elevated by difficulty, distinctness, possession, or purity.}
}
\]

\section*{28. Best Member and Best Kind}

If one kind is better than another without qualification, then the best member
of the first kind is better than the best member of the second kind.

\[
A > B \Rightarrow \text{best } A > \text{best } B
\]

Conversely, if the best member of one kind is better than the best member of
another, then the kind itself is better without qualification.

\[
\text{best } A > \text{best } B \Rightarrow A > B
\]

\subsection*{Example}

If man is better than horse, then the best man is better than the best horse.

If the best man is better than the best horse, then man is better than horse.

\subsection*{Teaching Point}

Compare kinds through their best instances, and instances through their kinds.

\[
\boxed{
\text{The best case can reveal the rank of the kind.}
}
\]

\section*{29. Goods Shareable with Friends}

Things our friends can share are more desirable than those they cannot.

\[
\text{shareable with friends} > \text{not shareable}
\]

Also, things we would rather really do for a friend are more desirable than
things we would rather merely seem to do for a stranger.

\subsection*{Example}

We would rather really do good to friends than merely seem to do so.

With strangers, we may care more about appearing to do good.

\subsection*{Teaching Point}

Friendship reveals which goods we value genuinely rather than merely
reputationally.

\[
\boxed{
\text{What we choose for friends reveals what we truly prize.}
}
\]

\section*{30. Superfluities and Necessities}

Superfluities are better than necessities in one respect, because the good life
is better than mere life.

\[
\text{good life} > \text{mere life}
\]

But necessities may be more desirable for someone who lacks them.

\subsection*{Example}

To be a philosopher is better than making money.

But for someone who lacks the necessities of life, making money may be more
desirable.

\subsection*{Teaching Point}

Better and more desirable can come apart under conditions of need.

\[
\boxed{
\text{Superfluities may be better, while necessities may be more urgent.}
}
\]

\section*{31. What Cannot Be Got from Another}

What cannot be obtained from another is more desirable than what can be
obtained from another.

\[
\text{not obtainable from another} > \text{obtainable from another}
\]

\subsection*{Example}

Justice is more desirable than courage in this respect.

A person cannot simply borrow another person's justice.

\subsection*{Teaching Point}

Goods that must be possessed personally have special priority.

\[
\boxed{
\text{A non-transferable good is especially choiceworthy.}
}
\]

\section*{32. Desirable Without Another}

A is more desirable if A is desirable without B, but B is not desirable without
A.

\[
A \text{ desirable without } B
\]

\[
B \text{ not desirable without } A
\]

\[
\therefore A > B
\]

\subsection*{Example}

Power is not desirable without prudence.

Prudence is desirable without power.

Therefore:

\[
\text{prudence} > \text{power}
\]

\subsection*{Teaching Point}

The good that makes another safe or choiceworthy outranks it.

\[
\boxed{
\text{A governing good outranks the power that needs it.}
}
\]

\section*{33. What We Wish to Be Thought to Possess}

If we repudiate one thing in order to be thought to possess another, then the
one we wish to be thought to possess is treated as more desirable.

\subsection*{Example}

Someone may repudiate love of hard work so that people will think he is a
genius.

\[
\text{repudiate hard work} \rightarrow \text{seem naturally gifted}
\]

This suggests that natural giftedness is treated as more desirable than mere
industry.

\subsection*{Teaching Point}

Reputational behavior reveals perceived rankings of goods.

\[
\boxed{
\text{What people want to seem to possess shows what they think is more admirable.}
}
\]

\section*{34. Blameworthiness of Vexation}

That is more desirable whose absence makes vexation less blameworthy.

That is also more desirable whose absence makes failure to be vexed more
blameworthy.

\subsection*{Simpler Form}

If losing something is a more legitimate cause of distress, that thing is more
desirable.

\[
\text{more justified grief at absence} \Rightarrow \text{more desirable good}
\]

\subsection*{Teaching Point}

Appropriate distress at loss can indicate the value of what is lost.

\[
\boxed{
\text{The nobler the loss, the more reasonable the grief.}
}
\]

\section*{35. Peculiar Virtue of the Species}

Of things belonging to the same species, the one that possesses the peculiar
virtue of the species is more desirable.

If both possess it, the one that possesses it in a greater degree is more
desirable.

\[
\text{same species}
\]

\[
\text{has proper excellence} > \text{lacks proper excellence}
\]

\subsection*{Teaching Point}

A thing is better when it better realizes the excellence proper to its kind.

\[
\boxed{
\text{The proper virtue of a kind supplies the standard for ranking its members.}
}
\]

\section*{36. What Makes Other Things Good}

If one thing makes good whatever it touches, while another does not, the former
is more desirable.

\[
\text{makes others good} > \text{does not}
\]

If both do so, the one that makes the better or more important object good is
more desirable.

\subsection*{Example}

\[
\text{makes the soul good} > \text{makes the body good}
\]

\subsection*{Teaching Point}

A good that communicates goodness has higher rank.

\[
\boxed{
\text{The power to improve what receives it is a mark of superior goodness.}
}
\]

\section*{37. Inflexions, Uses, Actions, and Works}

Aristotle says to judge things by their inflected forms, uses, actions, and
works, and to judge these by the things themselves.

\subsection*{Example}

If:

\[
\text{justly} > \text{courageously}
\]

then:

\[
\text{justice} > \text{courage}
\]

And if:

\[
\text{justice} > \text{courage}
\]

then:

\[
\text{justly} > \text{courageously}
\]

\subsection*{Teaching Point}

Value travels between a thing, its use, its action, its work, and its linguistic
relatives.

\[
\boxed{
\text{A good and its proper activity can be judged through one another.}
}
\]

\section*{38. Excess Over a Standard}

If one thing exceeds and another falls short of the same standard of good, the
one that exceeds is more desirable.

\[
A > \text{standard}
\]

\[
B < \text{standard}
\]

\[
\therefore A > B
\]

If both exceed the standard, the one that exceeds it more is more desirable.

\subsection*{Teaching Point}

Comparison can be made by measuring each item against a shared standard.

\[
\boxed{
\text{A shared standard makes comparative ranking clearer.}
}
\]

\section*{39. Excess of One Thing Over Another}

If the excess of one thing is more desirable than the excess of another, then
the thing itself is more desirable.

\subsection*{Example}

An excess of friendship is more desirable than an excess of money.

Therefore:

\[
\text{friendship} > \text{money}
\]

\subsection*{Teaching Point}

The value of excess can reveal the value of the thing exceeded.

\[
\boxed{
\text{If more of A is better than more of B, A is likely better than B.}
}
\]

\section*{40. What One Would Rather Have by One's Own Doing}

What a person would rather have by his own doing is more desirable than what he
would rather get by another's doing.

\subsection*{Example}

\[
\text{friends} > \text{money}
\]

A person would rather have friends through his own character and action, while
he might accept money from another.

\subsection*{Teaching Point}

Goods tied to one's own action and character have special desirability.

\[
\boxed{
\text{A self-achieved good may outrank an externally received good.}
}
\]

\section*{41. Addition}

Aristotle recommends judging by addition.

If adding A to the same thing as B makes the whole more desirable than adding B,
then A is more desirable than B.

\[
X + A > X + B
\]

\[
\therefore A > B
\]

\subsection*{Warning}

The shared context must be fair.

If the common term helps one added thing but not the other, the comparison is
misleading.

\subsection*{Example}

A saw may be more useful than a sickle when combined with carpentry, but that
does not prove that a saw is more desirable without qualification.

\subsection*{Teaching Point}

Addition reveals value only when the shared context is fair.

\[
\boxed{
\text{Do not choose a common context that secretly favors one side.}
}
\]

\section*{42. Subtraction}

Subtraction works similarly.

The thing whose removal leaves the lesser good may be taken as the greater
good.

\[
X - A < X - B
\]

\[
\therefore A > B
\]

\subsection*{Teaching Point}

The value of a thing may be revealed by what remains when it is removed.

\[
\boxed{
\text{The greater loss indicates the greater good.}
}
\]

\section*{43. Desired for Itself and Desired for Appearance}

A thing desired for itself is more desirable than what is desired merely for the
look of it.

\subsection*{Example}

\[
\text{health} > \text{beauty}
\]

Beauty may be desired for appearance.

Health is desirable even if no one sees it.

\subsection*{Definition}

Something is desired for the look of it if, supposing no one knew of it, one
would not care to have it.

\[
\text{desired only if seen} = \text{desired for appearance}
\]

\[
\text{desired even if unseen} = \text{desired for itself}
\]

\subsection*{Teaching Point}

What remains desirable without witnesses is more truly desirable.

\[
\boxed{
\text{The unseen test reveals whether a good is intrinsic or reputational.}
}
\]

\section*{44. Desired for Both Itself and Appearance}

If one thing is desirable both for itself and for appearance, while another is
desirable on only one ground, then the first is more desirable.

\[
\text{for itself + for appearance} > \text{for one ground only}
\]

\subsection*{Teaching Point}

A good supported by more kinds of desirability has comparative advantage.

\[
\boxed{
\text{Multiple grounds of desirability strengthen the case.}
}
\]

\section*{45. More Precious in Itself}

Whichever is more precious in itself is better and more desirable.

A thing is more precious in itself when we choose it for itself, without
expecting anything else to come from it.

\subsection*{Teaching Point}

Intrinsic preciousness is a strong sign of desirability.

\[
\boxed{
\text{What we choose even without further gain is more precious in itself.}
}
\]

\section*{46. Distinguish the Senses of Desirable}

Aristotle says we should distinguish the senses in which something is
desirable.

A thing may be desirable with a view to:

\[
\begin{array}{ll}
1. & \text{expediency} \\
2. & \text{honor} \\
3. & \text{pleasure}
\end{array}
\]

\subsection*{Teaching Point}

Clarify the kind of desirability before comparing desirables.

\[
\boxed{
\text{The question ``Which is more desirable?'' requires asking: desirable how?}
}
\]

\section*{47. Useful for More Ends}

What is useful for all or most ends is more desirable than what is useful in
fewer ways.

\[
\text{useful for honor, pleasure, and expediency}
>
\text{useful for only one}
\]

If the same characters belong to both, ask which possesses them more markedly.

\subsection*{Teaching Point}

Breadth and intensity of usefulness both affect desirability.

\[
\boxed{
\text{A good useful in more ways, or more strongly in the same ways, is preferable.}
}
\]

\section*{48. Serves the Better Purpose}

What serves the better purpose is more desirable.

\subsection*{Example}

\[
\text{what promotes virtue} > \text{what promotes pleasure}
\]

Likewise, what more obstructs the desirable is more objectionable.

\subsection*{Example}

Disease is more objectionable than ugliness because disease is a greater
hindrance to both pleasure and being good.

\subsection*{Teaching Point}

Means are ranked by the worth of the ends they serve or obstruct.

\[
\boxed{
\text{The higher the end, the higher the value of what serves it.}
}
\]

\section*{49. Mixed Desirability and Objectionability}

A thing that is both desirable and objectionable in like measure is less
desirable than something that is desirable only.

\[
\text{desirable + objectionable}
<
\text{desirable without comparable objection}
\]

\subsection*{Teaching Point}

A mixed good is lower than a good without a comparable drawback.

\[
\boxed{
\text{Desirability is weakened by equal objectionability.}
}
\]

\section*{50. From Comparative Rules to Simple Desirability}

The rules for comparing desirables can also show that something is simply
desirable or objectionable.

\subsection*{Example}

If what is more useful is more desirable, then what is useful is desirable.

\[
\text{more useful} \Rightarrow \text{more desirable}
\]

therefore:

\[
\text{useful} \Rightarrow \text{desirable}
\]

\subsection*{Teaching Point}

Comparative rules can yield non-comparative conclusions.

\[
\boxed{
\text{If more P implies more desirable, P itself may imply desirable.}
}
\]

\section*{51. Generalizing Comparative Rules}

Aristotle says that rules about comparative degrees and amounts should be taken
in the most general possible form.

The more general the rule, the more cases it can handle.

\subsection*{Example}

Specific rule:

\[
\text{What is good by nature is more desirable than what is not good by nature.}
\]

Generalized rule:

\[
\text{What naturally exhibits a quality exhibits it more than what does not naturally exhibit it.}
\]

\subsection*{Teaching Point}

A topical rule is most useful when expressed in its most general form.

\[
\boxed{
\text{Generalize the rule so that it applies beyond the first example.}
}
\]

\section*{52. Comparative Degree Rules in General Form}

A thing has a character in a greater degree if:

\begin{enumerate}[itemsep=0.25em]
    \item it has that character by nature,
    \item it imparts that character to what possesses it,
    \item it exceeds a shared standard while another falls short,
    \item it produces a greater result when added,
    \item its removal leaves the remainder with less of that character,
    \item it is freer from admixture with the contrary,
    \item it admits the definition of that character more fully.
\end{enumerate}

\subsection*{Example}

If the definition of white is:

\[
\text{a colour that pierces the vision}
\]

then what is more fully a colour that pierces the vision is whiter.

\subsection*{Teaching Point}

Comparative degree is often judged by nature, effect, standard, purity, and
definition.

\[
\boxed{
\text{To compare degrees, ask which case more fully realizes the character.}
}
\]

\section*{53. Particular Problems}

If the question is put in a particular rather than universal form, the universal
rules can still be used.

\[
\text{If true of all, true of some.}
\]

\[
\text{If untrue of all, untrue of some.}
\]

\subsection*{Teaching Point}

Universal constructive and destructive rules remain useful for particular
problems.

\[
\boxed{
\text{Universal proof can settle particular questions.}
}
\]

\section*{54. Useful Rules for Particular Problems}

For particular problems, Aristotle especially recommends rules drawn from:

\[
\begin{array}{ll}
1. & \text{opposites} \\
2. & \text{co-ordinates} \\
3. & \text{inflexions}
\end{array}
\]

\subsection*{Examples}

If some pleasure is good, then some pain may be evil.

If some pleasant thing is objectionable, then pleasure is in some case
objectionable.

If what happens justly is in some cases evil, then what happens unjustly may in
some cases be good.

\subsection*{Teaching Point}

Opposites and related forms allow particular claims to be transferred across
term-families.

\[
\boxed{
\text{Particular arguments often move through related terms and opposites.}
}
\]

\section*{55. Greater, Less, and Like Degrees in Particular Problems}

Greater, less, and like degrees can establish or overthrow particular claims.

\subsection*{More Likely Case}

If a character does not belong where it is more likely to belong, then it does
not belong where it is less likely.

\[
\text{not where more likely} \Rightarrow \text{not where less likely}
\]

\subsection*{Less Likely Case}

If a character belongs where it is less likely to belong, then it belongs where
it is more likely.

\[
\text{yes where less likely} \Rightarrow \text{yes where more likely}
\]

\subsection*{Teaching Point}

Comparative likelihood can settle particular questions.

\[
\boxed{
\text{Use stronger and weaker cases to test the case before you.}
}
\]

\section*{56. Hypothesis in Particular Problems}

Aristotle allows the use of hypothesis in particular problems.

One may claim that if an attribute belongs or does not belong in one case, it
does so in all similar cases.

\subsection*{Example}

If the soul of man is immortal, then other souls are immortal as well.

If this soul is not immortal, neither are the others.

\subsection*{Teaching Point}

Hypothesis can universalize a particular question.

\[
\boxed{
\text{A particular admission may be turned into a universal commitment.}
}
\]

\section*{57. Indefinite Statements}

An indefinite statement does not clearly specify whether it is universal or
particular.

\subsection*{Example}

\[
\text{Pleasure is good.}
\]

This may mean:

\[
\text{All pleasure is good.}
\]

or:

\[
\text{Some pleasure is good.}
\]

\subsection*{Teaching Point}

Indefinite claims are easier to establish than to demolish.

\[
\boxed{
\text{Ambiguous quantity gives the assertor flexibility.}
}
\]

\section*{58. Establishing an Indefinite Statement}

An indefinite affirmative can be established in two ways.

\[
\text{All pleasure is good.}
\]

or:

\[
\text{Some pleasure is good.}
\]

Either is enough to establish the indefinite claim:

\[
\text{Pleasure is good.}
\]

\subsection*{Teaching Point}

A weak particular proof may establish an indefinite assertion.

\[
\boxed{
\text{To establish the indefinite, one case may be enough.}
}
\]

\section*{59. Demolishing an Indefinite Statement}

An indefinite affirmative can be demolished only by proving the universal
opposite.

If someone says:

\[
\text{Pleasure is good.}
\]

and means some pleasure is good, it is not enough to show:

\[
\text{Some pleasure is not good.}
\]

To demolish the claim fully, one must show:

\[
\text{No pleasure is good.}
\]

\subsection*{Teaching Point}

To refute an indefinite claim, remove every possible interpretation that would
make it true.

\[
\boxed{
\text{Indefinite claims require universal refutation.}
}
\]

\section*{60. Definite Statements}

A definite statement gives a clearer target.

\subsection*{Example}

\[
\text{Only one pleasure is good.}
\]

This can be demolished in several ways:

\[
\text{All pleasure is good.}
\]

\[
\text{No pleasure is good.}
\]

\[
\text{More than one pleasure is good.}
\]

\subsection*{Teaching Point}

The more definite a claim is, the more exposed it is to refutation.

\[
\boxed{
\text{Precision increases vulnerability to counterargument.}
}
\]

\section*{61. Exclusive Claims}

If someone says:

\[
\text{Prudence alone among the virtues is knowledge,}
\]

there are four ways to demolish the claim.

\[
\text{All virtue is knowledge.}
\]

\[
\text{No virtue is knowledge.}
\]

\[
\text{Some other virtue, such as justice, is knowledge.}
\]

\[
\text{Prudence itself is not knowledge.}
\]

\subsection*{Teaching Point}

Exclusive claims can be refuted by attacking universality, exclusion, or the
favored case itself.

\[
\boxed{
\text{Only-claims expose several lines of attack.}
}
\]

\section*{62. Individual Instances and Division}

Aristotle returns to the method of individual instances and division by genera.

When an attribute is said to belong or not belong, examine individual cases.

Then divide genera into species until reaching the indivisible cases.

\[
\text{genus}
\rightarrow
\text{species}
\rightarrow
\text{individual cases}
\]

\subsection*{Teaching Point}

Even comparative dialectic returns to division, instances, and counterexamples.

\[
\boxed{
\text{The dialectician tests general claims by ordered descent into cases.}
}
\]

\section*{63. Forcing Universal Admission or Objection}

If, after division, the attribute appears to belong in all cases or no cases,
one may press the opponent.

The opponent must either:

\[
\text{admit the universal}
\]

or:

\[
\text{produce an objection}
\]

\subsection*{Teaching Point}

A well-conducted survey of cases forces either concession or counterexample.

\[
\boxed{
\text{Division turns scattered cases into dialectical pressure.}
}
\]

\section*{64. Definite Accidents by Species or Number}

When possible, make the accident definite either specifically or numerically.

\subsection*{Example: Time and Motion}

To show that time is not moved and is not a movement, enumerate the species of
movement.

If time belongs to none of them, then time is not moved and is not a movement.

\subsection*{Example: Soul and Number}

To show that the soul is not a number, divide all numbers into odd and even.

If the soul is neither odd nor even, then the soul is not a number.

\subsection*{Teaching Point}

A claim can be demolished by exhausting the species under the proposed genus.

\[
\boxed{
\text{If none of the species applies, the genus does not apply.}
}
\]

\section*{65. The Architecture of Book III}

Book III moves through the following arc:

\[
\text{comparative inquiry}
\rightarrow
\text{basic standards of choice}
\rightarrow
\text{ends vs. means}
\rightarrow
\text{in itself vs. accidental}
\]

\[
\rightarrow
\text{absolute vs. relative goods}
\rightarrow
\text{rank of subjects and causes}
\rightarrow
\text{consequences and number of goods}
\]

\[
\rightarrow
\text{timing and usefulness}
\rightarrow
\text{loss, acquisition, likeness}
\rightarrow
\text{friendship, necessity, prudence}
\]

\[
\rightarrow
\text{species excellence}
\rightarrow
\text{addition and subtraction}
\rightarrow
\text{kinds of desirability}
\]

\[
\rightarrow
\text{particular, indefinite, and definite problems}
\]

\subsection*{Teaching Point}

Book III teaches dialectical comparison of goods.

\[
\boxed{
\text{The question is not merely whether something is good, but how it ranks among goods.}
}
\]

\section*{66. Summary of Main Ideas}

\begin{enumerate}[itemsep=0.5em]
    \item Book III gives commonplace rules for comparing what is better or more
    desirable.

    \item Comparative dialectic is most useful when the alternatives are close
    enough to be genuinely debated.

    \item A single relevant advantage may decide a close comparison.

    \item More lasting and secure goods are more desirable.

    \item The choices of the prudent, good, lawful, or expert supply standards
    for comparison.

    \item What is desired for itself is more desirable than what is desired for
    something else.

    \item What is desirable in itself outranks what is desirable per accidens.

    \item What is good absolutely outranks what is good only under special
    conditions.

    \item Ends outrank means, and nearer means outrank more remote means.

    \item The rank of an end can raise the rank of what produces it.

    \item Noble, precious, and praiseworthy things are more desirable than
    merely useful things.

    \item Consequences reveal hidden superiority when alternatives are close.

    \item More goods are better than fewer goods only when the added goods have
    independent value.

    \item Pleasure and pain modify desirability.

    \item Timing and season affect comparative desirability.

    \item Constant usefulness is a mark of greater desirability.

    \item Loss, destruction, acquisition, and generation reveal value.

    \item Likeness to the good is useful only when the resemblance is relevant.

    \item Better and more desirable can come apart under conditions of need.

    \item Prudence outranks power because power is not desirable without
    prudence.

    \item The proper excellence of a species supplies a standard for comparing
    members of that species.

    \item Addition and subtraction can reveal comparative value when the context
    is fair.

    \item What remains desirable without witnesses is more truly desirable than
    what is desired only for appearance.

    \item Desirability must be distinguished according to expediency, honor, and
    pleasure.

    \item Comparative rules can yield simple conclusions about desirability.

    \item Indefinite claims are easier to establish than to demolish.

    \item Definite and exclusive claims are more vulnerable to refutation.

    \item Division by genus and species remains a central tool of dialectic.
\end{enumerate}

\section*{Final Framing}

\[
\boxed{
\textit{Topics}, \text{ Book III, teaches dialectical comparison of goods.}
}
\]

Book II teaches how to test whether a predicate belongs accidentally. Book III
teaches how to compare goods, ends, means, consequences, losses, acquisitions,
and circumstances when the question is which thing is more desirable.

\[
\boxed{
\text{Dialectic does not only ask what is true; it also asks what is worth choosing.}
}
\]

\end{document}